\section{格点矩阵元}\label{sec:lattice-details}

对于 K 介子基态 $|0 (P_z) \rangle$，我们在若干助推动量 $P_z$ 与 Wilson 线位移长度 $z$ 下计算其格点胶子矩阵元，所用非极化胶子算符最早由文献~\cite{Balitsky:2019krf} 引入：
%
\begin{equation}\label{eq:gluon_operator}
 {\cal O}(z)\equiv\sum_{i\neq z,t}{\cal O}(F^{ti},F^{ti};z)-\sum_{i,j\neq z,t}{\cal O}(F^{ij},F^{ij};z),
\end{equation}
%
其中 ${\cal O}(F^{\mu\nu}, F^{\alpha\beta};z) = F^\mu_\nu(z)U(z,0)F^\alpha_\beta(0)$，Wilson 链长度记为 $z$，场强 $F_{\mu\nu}$ 由下式给出：
%
\begin{equation}\label{field_strength}
 F_{\mu\nu} = \frac{i}{8a^2g_0}\left(\mathcal{P}_{[\mu,\nu]}+\mathcal{P}_{[\nu,-\mu]}+\mathcal{P}_{[-\mu,-\nu]}+\mathcal{P}_{[-\nu,\mu]}\right),
\end{equation}
%
其中 $a$ 是格距，$g_0$ 是强耦合常数；闭合方格（plaquette）$\mathcal{P_{\mu,\nu}}=U_\mu(x)U_\nu(x+a\hat{\mu})U^\dag_\mu(x+a\hat{\nu})U^\dag_\nu(x)$，且 $\mathcal{P_{[\mu,\nu]}}=\mathcal{P}_{\mu,\nu}-\mathcal{P}_{\nu,\mu}$。
算符 $\sum_{i\neq z,t}{\cal O}(F^{ti},F^{zi};z)$ 对应于与文献~\cite{Balitsky:2019krf} 相同的匹配核。
然而该算符并未被采用，因为由于运动学原因它在 $P_z=0$ 处为零，这会增加获得分布的难度。
在计算出上述矩阵元之后，我们考察它们对 Ioffe 时间 $\nu=zP_z$ 的依赖性。

在格点上，我们在三个系综上采用 clover 价费米子；这些系综含 $N_f=2+1+1$ 味高度改进交错夸克（HISQ）~\cite{Follana:2006rc}，由 MILC 合作组生成~\cite{MILC:2010pul,Bazavov:2012xda}，涵盖两个格距（$a\approx 0.12$ 与 0.15~fm）与 310~MeV 的 $\pi$ 介子质量。
clover 夸克质量经过调节，以重现 PNDME 合作组所用的最轻奇异及轻海赝标介子~\cite{Rajan:2017lxk,Bhattacharya:2015wna,Bhattacharya:2015esa,Bhattacharya:2013ehc}。
如文献~\cite{Fan:2018dxu} 所阐明，我们对胶子圈施加五步 HYP 光滑化（smearing）~\cite{Hasenfratz:2001hp} 以降低统计噪声。
为达到更高的介子助推动量，我们对夸克场使用高斯型动量涂抹~\cite{Bali:2016lva}。
表~\ref{table-data} 给出了三个系综的格距 $a$、价 $\pi$ 介子质量 $M_\pi^\text{val}$ 与 K 介子质量 $M_K^\text{val}$、格点尺寸 $L^3\times T$、规范组态数目 $N_\text{cfg}$，以及两点关联函数的总测量次数 $N_\text{meas}^\text{2pt}$。


\begin{table}[!htbp]
\centering
\begin{tabular}{|c|c|c|}
\hline
  系综 &  a12m310 & a15m310 \\
\hline
  $a$ (fm)   & $0.1207(11)$  & $0.1510(20)$ \\
\hline
  $M_\pi^\text{val}$ (MeV)  & $311.1(6)$ & $319.1(31)$\\
\hline
  $M_{k}^\text{val}$ (MeV)  & 528.0(5) &  531.7(23)  \\
\hline
  $L^3\times T$   & $24^3\times 64$ & $16^3\times 48$ \\
\hline
  $P_z$ (GeV)   &  $[0,2.14]$ &  $[0,2.05]$ \\
\hline
  $N_\text{cfg}$  & 1013 & 900 \\
\hline
  $N_\text{meas}^\text{2pt}$   & 324,160 & 21,600 \\
\hline
%\hline
\end{tabular}
\caption{
本研究分析的、由 MILC 合作组生成的 HISQ 系综（配以 $N_f=2+1+1$ clover 价费米子用于三点关联函数拟合）的格距 $a$、价 $\pi$ 介子质量 $M_\pi^\text{val}$ 与 K 介子质量 $M_{k}^\text{val}$、格点尺寸 $L^3\times T$、规范组态数目 $N_\text{cfg}$、两点关联函数总测量次数 $N_\text{meas}^\text{2pt}$ 以及分离时间 $t_\text{sep}$。
}\label{table-data}
\end{table}




对一个介子 $\Phi$，两点关联函数由下式给出：
%
\begin{align}
C_\Phi^\text{2pt}(P_z;t)&=
 \int d^3y\, e^{-i y\cdot P_z} \langle \chi_\Phi(\vec{y},t)|\chi_\Phi(\vec{0},0)\rangle \nonumber \\
  &= |A_{\Phi,0}|^2 e^{-E_{\Phi,0}t} + |A_{\Phi,1}|^2 e^{-E_{\Phi,1}t} + ...,
\label{eq:2pt_fit_formula}
\end{align}
%
其中 $P_z$ 是介子沿 $z$ 方向的动量，$\chi_\Phi=\bar{q}_1\gamma_5 q_2$ 是赝标介子插值算符，$t$ 是欧氏时间。
$|A_{\Phi,i}|^2$ 与 $E_{\Phi,i}$ 分别表示基态（$i=0$）和第一激发态（$i=1$）的振幅与能量。
我们将 K 介子两点关联函数与胶子算符相乘以生成三点胶子关联函数。
随后把三点关联函数拟合到能量本征态展开的前几项，从而提取胶子算符的矩阵元：
%
\begin{align}\label{eq:3ptC}
&C_\Phi^\text{3pt}(z,P_z;t_\text{sep},t) \nonumber \\
&=\int d^3y\, e^{-iy\cdot P_z}\langle \chi_\Phi(\vec y,t_\text{sep})|{\cal O}(z,t)|\chi_\Phi(\vec 0,0)\rangle \nonumber \\
&= |A_{\Phi,0}|^2\langle 0|{\cal O}|0\rangle e^{-E_{\Phi,0}t_\text{sep}} \nonumber \\
&+ |A_{\Phi,0}||A_{\Phi,1}|\langle 0|{\cal O}|1\rangle e^{-E_{\Phi,1}(t_\text{sep}-t)}e^{-E_{\Phi,0}t} \nonumber \\
&+ |A_{\Phi,0}||A_{\Phi,1}|\langle 1|{\cal O}|0\rangle e^{-E_{\Phi,0}(t_\text{sep}-t)}e^{-E_{\Phi,1}t} \nonumber \\
&+ |A_{\Phi,1}|^2\langle 1|{\cal O}|1\rangle e^{-E_{\Phi,1}t_\text{sep}}+ ...,
\end{align}
%
其中 $t_{\text{sep}}$ 是源—汇时间分离，$t$ 是胶子算符的插入时间。
振幅 $A_{\Phi,0}$、$A_{\Phi,1}$ 与能量 $E_{\Phi,0}$、$E_{\Phi,1}$ 通过对两点关联函数的两态拟合提取。
$\langle 0|{\cal O}|0\rangle$、$\langle 0|{\cal O}|1\rangle$（$\langle 1|{\cal O}|0\rangle$）与 $\langle 1|{\cal O}|1\rangle$ 分别表示基态矩阵元、基态—激发态矩阵元与激发态矩阵元。




基态矩阵元是通过对三点关联函数做两态拟合，或跨多个时间分离做两态同步（two-sim）拟合获得的。
为验证拟合出的矩阵元被正确提取，我们计算了比值
%
\begin{equation}\label{eq:Ratio}
R^\text{ratio}(z,P_z;t_\text{sep},t) = \frac{C^\text{3pt}(z,P_z;t_\text{sep},t)}{C^\text{2pt}(P_z;t_\text{sep})};
\end{equation}
%
即三点与两点关联函数之比。
图~\ref{fig:a12m310_Ratio_figure} 与图~\ref{fig:a15m310_Ratio_figure} 的第一列展示了系综 a12m310 和 a15m310 在选定动量 $P_z$ 与 Wilson 线长度 $z$ 下 K 介子的 $R^\text{ratio}$。
若不存在激发态，式~\ref{eq:Ratio} 将直接给出基态矩阵元。
可以看到，$R^\text{ratio}$ 随源—汇分离的增大而增大。
在某些情形下，$R^\text{ratio}$ 在大分离处开始收敛，表明来自激发态的贡献逐渐减小。
灰色带描绘了在 $t_\text{sep}\in [5,9]$ 下对三点关联函数做“two-sim”拟合得到的基态矩阵元。
在第二列与第三列中，可以看出拟合基态矩阵元随 $t_\text{sep}^\text{min}$ 与 $t_\text{sep}^\text{max}$ 不同取值的变化。
在图~\ref{fig:a15m310_Ratio_figure} 所示的例子中，无论 $t_\text{sep}^{\text{min},\text{max}}$ 如何选取，a15m310 系综的基态拟合矩阵元都是稳定的。
然而，对于来自 a12m310 系综的例子，我们发现基态矩阵元对 $t_\text{sep}$ 的选取较为敏感。
分析中所采用的 $t_\text{sep}\in [5,9]$ 确实能可靠地提取基态。



\begin{figure*}[!tbp]
  \centering
\includegraphics[width=0.35\textwidth]{images/a12m310_ratio_fit_kaon_p1z4.pdf}
\includegraphics[width=0.2\textwidth]{images/a12m310_twosimfit_kaon_p1z4_tsmin.pdf}
\includegraphics[width=0.2\textwidth]{images/a12m310_twosimfit_kaon_p1z4_tsax.pdf}\\
\includegraphics[width=0.35\textwidth]{images/a12m310_ratio_fit_kaon_p2z1.pdf}
\includegraphics[width=0.2\textwidth]{images/a12m310_twosimfit_kaon_p2z1_tsmin.pdf}
\includegraphics[width=0.2\textwidth]{images/a12m310_twosimfit_kaon_p2z1_tsmax.pdf}
\caption{$a\approx 0.12$~fm、$M_\pi\approx 310$~MeV 的 $\pi$ 介子质量系综的示例比值图（左列）
与 two-sim 拟合（后两列），对应 $P_z=8\pi/L$、$z=4$（上排）与 $P_z=4\pi/L$、$z=1$（下排）。
所有图中灰色带均为采用 $t_\text{sep}\in[5,9]$ 的 two-sim 拟合提取的基态矩阵元。
从左到右各列分别为：
三点与两点关联函数之比（附采用 $t_\text{sep}\in [5,9]$ 的 two-sim 拟合重建带），表示为 $t-t_\text{sep}/2$ 的函数；
采用 $t_\text{sep}\in[t_\text{sep}^\text{min},9]$ 的 two-sim 拟合结果随 $t_\text{sep}^\text{min}$ 的变化；
以及采用 $t_\text{sep}\in[5,t_\text{sep}^\text{max}]$ 的 two-sim 拟合结果随 $t_\text{sep}^\text{max}$ 的变化。
}
\label{fig:a12m310_Ratio_figure}
\end{figure*}




\begin{figure*}[!tbp]
  \centering
\includegraphics[width=0.35\textwidth]{images/a15m310_ratio_fit_kaon_p1z4.pdf}
\includegraphics[width=0.2\textwidth]{images/a15m310_twosimfit_kaon_p1z4_tsmin.pdf}
\includegraphics[width=0.2\textwidth]{images/a15m310_twosimfit_kaon_p1z4_tsax.pdf}\\
\includegraphics[width=0.35\textwidth]{images/a15m310_ratio_fit_kaon_p2z1.pdf}
\includegraphics[width=0.2\textwidth]{images/a15m310_twosimfit_kaon_p2z1_tsmin.pdf}
\includegraphics[width=0.2\textwidth]{images/a15m310_twosimfit_kaon_p2z1_tsmax.pdf}
\caption{$a\approx 0.15$~fm、$M_\pi\approx 310$~MeV 的 $\pi$ 介子质量系综的示例比值图（左列）
与 two-sim 拟合（后两列），对应 $P_z=1\times 2\pi/L$、$z=1$（上排）与 $P_z=4\times 2\pi/L$、$z=4$（下排）。
所有图中灰色带均为采用 $t_\text{sep}\in[5,9]$ 的 two-sim 拟合提取的基态矩阵元。
从左到右各列分别为：
三点与两点关联函数之比（附采用 $t_\text{sep}\in [5,9]$ 的 two-sim 拟合重建带），表示为 $t-t_\text{sep}/2$ 的函数；
采用 $t_\text{sep}\in[t_\text{sep}^\text{min},9]$ 的 two-sim 拟合结果随 $t_\text{sep}^\text{min}$ 的变化；
以及采用 $t_\text{sep}\in[5,t_\text{sep}^\text{max}]$ 的 two-sim 拟合结果随 $t_\text{sep}^\text{max}$ 的变化。
}
\label{fig:a15m310_Ratio_figure}
\end{figure*}
